
\medskip
\phantomsection\label{ex:sga3-VIB-6.2.4}%
\noindent\textbf{Examples 6.2.4.---}
Let us give important examples of functors which reduce to
functors
\[
  \prod_{Z/S}Y/Z
\]
of the type considered in
\hyperref[thm:sga3-VIB-6.2.3]{6.2.3}, and for which it will be
useful to have representability criteria later on. We denote by
\(S\) a scheme, and by \(X,Y,Z,\ldots\) schemes over \(S\).

\medskip
\phantomsection\label{ex:sga3-VIB-6.2.4-a}%
\noindent\textbf{a)}
Suppose given an \(S\)-morphism
\begin{equation}
  q:X\longrightarrow\underline{\operatorname{Hom}}_S(Y,Z),
  \tag{x}\label{eq:sga3-VIB-6.2.4-x}
\end{equation}
(``\(X\) acts on \(Y\), with values in \(Z\)''), that is, a
morphism
\begin{equation}
  r:X\times_S Y\longrightarrow Z.
  \tag{xx}\label{eq:sga3-VIB-6.2.4-xx}
\end{equation}
Consider a subscheme \(Z'\) of \(Z\), and hence a monomorphism
\[
  \underline{\operatorname{Hom}}_S(Y,Z')
  \longrightarrow
  \underline{\operatorname{Hom}}_S(Y,Z),
\]
which makes the first functor a subfunctor of the second. Let
\(X'\) be the inverse image of this subfunctor under \(q\). Thus
\(X'\) is the subfunctor of \(X\) for which \(X'(T)\) is the set
of \(x\in X(T)\) such that
\[
  q(x):Y_T\longrightarrow Z_T
\]
factors through \(Z'_T\).

This functor \(X'\) may be described as follows. Put
\[
  P=X\times_S Y,
\]
and let \(P'\) be the inverse image of \(Z'\) under
\(r:P\to Z\). Then there is an evident isomorphism
\begin{equation}
  X'\simeq\prod_{P/X}P'/P.
  \tag{xxx}\label{eq:sga3-VIB-6.2.4-xxx}
\end{equation}
It follows that, if \(Y\) is essentially free over \(S\) and
\(Z'\) is closed in \(Z\), the subfunctor \(X'\) of \(X\) is
representable by a closed subscheme of \(X\).%
\footnote{\label{fnref:VI_B-56}\label{note:sga3-VIB-56}Editor's note: we have added the
sentence that follows.}
If, moreover, \(Z'\to Z\) is of finite presentation, then so is
\(X'\to X\).

\medskip
\phantomsection\label{ex:sga3-VIB-6.2.4-b}%
\noindent\textbf{b)}
Suppose given two ways of letting \(X\) act on \(Y\), with values
in \(Z\), that is, two morphisms
\[
  q_1,q_2:X\rightrightarrows
  \underline{\operatorname{Hom}}_S(Y,Z),
\]
and put
\[
  X'=\operatorname{Ker}(q_1,q_2).
\]
Thus \(X'\) is the subfunctor of \(X\) for which \(X'(T)\) is the
set of \(x\in X(T)\) such that the two morphisms
\[
  q_1(x),q_2(x):Y_T\rightrightarrows Z_T
\]
are equal. The data of \(q_1,q_2\) are equivalent to the data of
a morphism
\[
  q:X\longrightarrow
  \underline{\operatorname{Hom}}_S(Y,Z\times_S Z),
\]
or, equivalently, of a morphism
\[
  r:X\times_S Y\longrightarrow Z\times_S Z.
\]
Put \(U=Z\times_S Z\), and let \(U'\) be the diagonal subscheme
of \(Z\times_S Z\). Then \(X'\) is precisely the inverse image
under \(q\) of the subfunctor
\[
  \underline{\operatorname{Hom}}_S(Y,U')
  \subset
  \underline{\operatorname{Hom}}_S(Y,U).
\]
It can therefore be put in the form
\hyperref[eq:sga3-VIB-6.2.4-xxx]{\textup{(xxx)}}, with
\(P=X\times_S Y\) and \(P'\) the inverse image of the diagonal
under \(r\), that is, the kernel of
\[
  r_1,r_2:X\times_S Y\rightrightarrows Z.
\]
We are thus under the hypotheses of
\hyperref[ex:sga3-VIB-6.2.4-a]{\textup{(a)}}.

Consequently, if \(Y\) is essentially free over \(S\) and \(Z\)
is separated over \(S\), the subfunctor \(X'\) of \(X\) is
representable by a closed subscheme of \(X\).%
\textsuperscript{\normalfont\ref{fnref:VI_B-56}}
If, moreover, \(Z\to S\) is locally of finite type, then
\(X'\to X\) is of finite presentation.

\medskip
\phantomsection\label{ex:sga3-VIB-6.2.4-c}%
\noindent\textbf{c)}
Suppose given a morphism
\[
  q:X\longrightarrow\underline{\operatorname{Hom}}_S(Y,Y),
\]
that is, ``\(X\) acts on \(Y\).'' Let \(X'\) be the ``kernel'' of
this morphism: \(X'\) is the subfunctor of \(X\) for which
\(X'(T)\) is the set of \(x\in X(T)\) such that
\[
  q(x):Y_T\longrightarrow Y_T
\]
is the identity. This functor falls under
\hyperref[ex:sga3-VIB-6.2.4-b]{\textup{(b)}}, as one sees by
introducing a second morphism
\[
  q':X\longrightarrow\underline{\operatorname{Hom}}_S(Y,Y)
\]
by letting \(X\) act trivially on \(Y\). Hence, if \(Y\) is
essentially free and separated over \(S\), the subfunctor \(X'\),
the kernel of \(q\), is representable by a closed subscheme of
\(X\).%
\textsuperscript{\normalfont\ref{fnref:VI_B-56}}
If, moreover, \(Y\to S\) is locally of finite type, then
\(X'\to X\) is of finite presentation.

\medskip
\phantomsection\label{ex:sga3-VIB-6.2.4-d}%
\noindent\textbf{d)}
Under the hypotheses of \hyperref[ex:sga3-VIB-6.2.4-c]{(c)},
consider the subfunctor \(Y'\) of \(Y\) ``of the invariants under
\(X\).'' Thus \(Y'(T)\) is the set of \(y\in Y(T)\) such that
the corresponding morphism
\[
  \bar q(y):X_T\longrightarrow Y_T
\]
is the constant \(T\)-morphism with value \(y\). Introducing
\(\bar q'\) as in \hyperref[ex:sga3-VIB-6.2.4-c]{(c)}, and the
morphisms corresponding to \(\bar q\) and \(\bar q'\),
\[
  \bar q,\bar q':Y\rightrightarrows
  \underline{\operatorname{Hom}}_S(X,Y),
\]
one sees that
\[
  Y'=\operatorname{Ker}(\bar q,\bar q').
\]
It therefore falls again under
\hyperref[ex:sga3-VIB-6.2.4-b]{\textup{(b)}}, with the roles of
\(X\) and \(Y\) reversed and with \(Z=Y\).

Consequently, if \(X\) is essentially free over \(S\) and \(Y\)
is separated over \(S\), the subfunctor \(Y'\) of \(Y\) of the
invariants under \(X\) is representable by a closed subscheme of
\(Y\).%
\textsuperscript{\normalfont\ref{fnref:VI_B-56}}
If, moreover, \(Y\to S\) is locally of finite type, then
\(Y'\to Y\) is of finite presentation.

\medskip
\phantomsection\label{ex:sga3-VIB-6.2.4-e}%
\noindent\textbf{e)}
Constructions of the kind described in the preceding examples
occur especially often in group theory. Thus, when \(G\) is an
\(S\)-group scheme acting on the \(S\)-scheme \(X\),
\[
  q:G\longrightarrow\underline{\operatorname{Aut}}_S(X),
\]
the kernel of \(q\), the subgroup of \(G\) acting trivially, is a
closed subscheme of \(G\), provided \(X\) is essentially free and
separated over \(S\)
(\hyperref[ex:sga3-VIB-6.2.4-c]{example~(c)}). The subobject
\(X^G\) of invariants is a closed subscheme of \(X\), provided
\(G\) is essentially free over \(S\) and \(X\) is separated over
\(S\) (\hyperref[ex:sga3-VIB-6.2.4-d]{example~(d)}).

Let \(Y,Z\) be subschemes of \(X\), and consider the subfunctor
\(\underline{\operatorname{Transp}}_G(Y,Z)\) of \(G\), the
transporter of \(Y\) into \(Z\). Its \(T\)-valued points, for
\(T\) over \(S\), are the elements \(g\in G(T)\) such that the
corresponding automorphism of \(X_T\) satisfies
\[
  g(Y_T)\subset Z_T,
\]
that is, induces a morphism \(Y_T\to X_T\) which factors through
\(Y_T\to Z_T\). Hence, if \(Y\) is essentially free over \(S\)
and \(Z\) is closed in \(X\), then
\(\underline{\operatorname{Transp}}_G(Y,Z)\) is a closed
subscheme of \(G\)
(\hyperref[ex:sga3-VIB-6.2.4-a]{example~(a)}).

One may also consider the strict transporter of \(Y\) into
\(Z\),%
whose \(T\)-valued points, for \(T\) over \(S\), are the elements
\(g\in G(T)\) such that \(g(Y_T)=Z_T\). This functor is precisely
\[
  \underline{\operatorname{Transp}}_G(Y,Z)
  \cap
  \sigma\bigl(\underline{\operatorname{Transp}}_G(Z,Y)\bigr),
\]
where \(\sigma\) is the inversion morphism of \(G\). Consequently,
if \(Y\) and \(Z\) are essentially free over \(S\) and closed in
\(X\), the strict transporter of \(Y\) into \(Z\) is a closed
subscheme of \(G\).

An important case is \(X=G\), with \(G\) acting on itself by
inner automorphisms. If \(H\) is a subscheme of \(G\), the strict
transporter of \(H\) into \(H\) is also called the
\emph{normalizer of \(H\) in \(G\)}, and is denoted
\(\underline{\operatorname{Norm}}_G H\). Hence, if \(H\) is a
closed group subscheme of \(G\), essentially free over \(S\),
then \(\underline{\operatorname{Norm}}_G H\) is representable by
a closed group subscheme of \(G\).

Finally, let \(Z\) be a subscheme of \(G\). Its centralizer
\(\underline{\operatorname{Centr}}_G(Z)\) in \(G\) is the group
subfunctor of \(G\) defined by the procedure of
\hyperref[ex:sga3-VIB-6.2.4-d]{(d)}, when one regards \(Z\) as
acting on \(G\) by the operations induced by those of \(G\).
Thus, if \(Z\) is essentially free over \(S\) and \(G\) is
separated over \(S\), then
\(\underline{\operatorname{Centr}}_G(Z)\) is a closed group
subscheme of \(G\). In particular, if \(G\) is essentially free
and separated over \(S\), then the center \(C\) of \(G\), which
is none other than \(\underline{\operatorname{Centr}}_G(G)\), is
a closed group subscheme of \(G\).

When \(S\) is the spectrum of a field,
\hyperref[prop:sga3-VIB-6.2.2]{6.2.2\textup{(b)}} shows that the
``essentially free'' conditions in examples
\hyperref[ex:sga3-VIB-6.2.4-a]{(a)}--%
\hyperref[ex:sga3-VIB-6.2.4-e]{(e)} above are automatically
satisfied, leaving only the separation conditions. Recalling
that a group scheme over a field is necessarily separated
(\hyperref[sga3:VIA:0.3]{VI\(_A\), 0.3}), we obtain, for example:

\medskip
\phantomsection\label{cor:sga3-VIB-6.2.5}%
\noindent\textbf{Corollary 6.2.5.---}
\emph{Let \(G\) be a group scheme over a field \(k\), and let
\(Y,Y'\) be two subschemes of \(G\). Then:}
\begin{enumerate}
\item[\textup{(i)}]\phantomsection\label{concl:sga3-VIB-6.2.5-i}
  \emph{The centralizer of \(Y\) in \(G\) is a closed group
  subscheme of \(G\); in particular, this holds for the center
  \(\underline{\operatorname{Centr}}_G(G)\) of \(G\).}
\item[\textup{(i\(^{\prime}\))}]
  \phantomsection\label{concl:sga3-VIB-6.2.5-iprime}
  \emph{More generally, if \(u,v:X\to G\) are morphisms of
  schemes, then \(\underline{\operatorname{Transp}}_G(u,v)\) is
  representable by a closed subscheme of \(G\).}
\item[\textup{(ii)}]\phantomsection\label{concl:sga3-VIB-6.2.5-ii}
  \emph{If \(Y\) is closed, the transporter
  \(\underline{\operatorname{Transp}}_G(Y',Y)\) is a closed
  subscheme of \(G\). If \(Y'\) is also closed, the same
  conclusion holds for
  \(\underline{\operatorname{Transp}}^{\mathrm{str}}_G(Y',Y)\).}
\item[\textup{(iii)}]\phantomsection\label{concl:sga3-VIB-6.2.5-iii}
  \emph{For every group subscheme}%
  \emph{\(H\) of \(G\),
  \(\underline{\operatorname{Norm}}_G(H)\) is a closed group
  subscheme of \(G\).}
\end{enumerate}
